
\section*{Техническое задание}

\begin{enumerate}
    \item Основание для выполнения работы~--- учебный план СПбПУ ИКНК 02.03.01.
    \item Исполнитель работы~--- студент группы 5130201/40002 Семенов Илья Андреевич.
    \item Цели, задачи и исходные данные.
        \begin{enumerate}[label=\alph*.]
            \item Целью работы является разработка программного комплекса на языке Haskell для статистического анализа растровых шрифтов в формате \verb|BDF|, позволяющего выявлять глифы с аномальными метриками и автоматически подбирать подходящую замену из встроенной базы эталонных глифов.
            \item В процессе выполнения работы требуется решить следующие задачи.
            \begin{enumerate}
                \item Изучить спецификацию формата растровых шрифтов BDF.
                \item Разработать парсер файлов формата BDF на языке Haskell.
                \item Реализовать математические модели для оценки метрик глифов (читаемость, пропорциональность, плотность заполнения).
                \item Разработать алгоритм обнаружения глифов с аномальными метриками.
                \item Разработать алгоритм поиска визуально похожих глифов и подбора замены из базы эталонов.
                \item Сформировать встроенную базу эталонных глифов для типовых символов.
                \item Реализовать консольный интерфейс для взаимодействия с пользователем.
                \item Обеспечить логирование действий и обработку ошибок ввода.
                \item Покрыть ключевые модули тестами с использованием библиотеки QuickCheck.
                \item Подготовить отчётную документацию.
            \end{enumerate}
            \item Исходными данными для проведения работы являются файлы растровых шрифтов в формате \verb|.bdf|.
        \end{enumerate}
    \item Основное содержание работы.
    \begin{enumerate}[label=\alph*.]
        \item Разработка проекта.
        \begin{enumerate}
            \item Проектирование типов данных для представления шрифта и глифов.
            \item Реализация парсера файла шрифта формата \verb|.bdf|.
            \item Реализация модулей анализа шрифта и расчёта метрик глифов.
            \item Реализация модуля подбора замены глифов из базы эталонов.
            \item Интеграция всех модулей в единый программный комплекс.
            \item Покрытие программного комплекса тестами.
        \end{enumerate}
        \item Отчётная документация.
        \begin{enumerate}
            \item Пояснительная записка с описанием архитектуры, алгоритмов и результатов тестирования.
            \item Руководство пользователя.
        \end{enumerate}
    \end{enumerate}
    \item Основные требования к выполнению работы.
    \begin{enumerate}[label=\alph*.]
        \item Разрабатываемый программный комплекс предназначен для автоматической валидации растровых шрифтов, используемых во встраиваемых системах\footnote{Встраиваемая система~--- это вычислительное устройство, встроенное в состав другого технического изделия и предназначенное для управления некоторыми процессами этого изделия.} (счётчики, терминалы оплаты, бытовая электроника) перед компиляцией прошивки, с целью предотвращения ошибок отображения на устройствах с ограниченными ресурсами и низким разрешением дисплеев.
        \item Состав разрабатываемого программного комплекса.
        \begin{enumerate}
            \item Модуль ввода-вывода: парсинг шрифтовых файлов \verb|.bdf| и запись отчёта в файл и консоль.
            \item Модуль расчёта метрик глифов и анализа шрифта.
            \item Модуль подбора замены глифов из базы эталонов.
            \item Модуль взаимодействия с пользователем~--- консольный интерфейс.
        \end{enumerate}
        \item Программное обеспечение должно обеспечить следующие
        технические характеристики по назначению.
        \begin{enumerate}
            \item Язык реализации~--- Haskell 2010.
            \item Система сборки~--- Stack.
            \item Запуск в ОС Windows $\geq 10$, Linux и macOS.
            \item Отсутствие внешних зависимостей, требующих компиляции C-кода.
        \end{enumerate}
        \item Функциональные требования.
        \begin{enumerate}
            \item Чтение и разбор файлов формата \verb|.bdf|.
            \item Расчёт и анализ метрик глифов шрифта.
            \item Обнаружение глифов с аномальными метриками и подбор замены из базы эталонов.
            \item Запись отчёта о результатах анализа в файл.
            \item Взаимодействие с пользователем через консольный интерфейс (меню действий и выбор параметров для анализа).
            \item Логирование действий пользователя.
            \item Обработка ошибок: корректная реакция на повреждённые файлы (или отсутствие файлов) и некорректный ввод пользователя.
        \end{enumerate}
        \item При разработке математического, программного и информационного
        обеспечения должны преимущественно использоваться типовые и
        стандартные средства, существующие библиотечные решения.
    \end{enumerate}
    \item Этапы и сроки выполнения работы.
    \begin{enumerate}[label=\alph*.]
        \item Этап~1~--- утверждённое техническое задание для курсовой работы. Срок: до 17.04.2026.
        \item Этап~2~--- созданная диаграмма типов и модулей. Срок: до 24.04.2026.
        \item Этап~3~--- реализованный проект на Stack на языке программирования Haskell. Срок: до 15.05.2026.
        \item Этап~4~--- защищённый отчёт по курсовой работе. Срок: до 22.05.2026.
    \end{enumerate}
    \item Результаты работы.
    \begin{enumerate}[label=\alph*.]
        \item Этап~1~--- документ с техническим заданием.
        \item Этап~2~--- диаграмма модулей с описанием типов данных и их взаимодействий.
        \item Этап~3~--- работающий программный комплекс на языке Haskell с тестами.
        \item Этап~4~--- отчёт по курсовой работе.
    \end{enumerate}
\end{enumerate}
